% =============================================================================
% Kapitel 1 -- Einleitung
% =============================================================================
\section{Einleitung}

\subsection{Motivation}

Entwurfsmuster sind bewährte Lösungsschablonen für wiederkehrende Probleme in der
objektorientierten Softwareentwicklung. Das Standardwerk der \textit{Gang of Four}
(Gamma et al., 1994) hat sie als gemeinsame Sprache für Architekturentscheidungen
etabliert. Ihr eigentlicher Wert zeigt sich jedoch erst im Kontext eines realen
Entwurfsproblems, die reine Theorie lässt ihren Nutzen oft abstrakt wirken.

Dieses Projekt setzt neun GoF-Patterns in \textbf{CardSiege} um, einem
rundenbasierten Kartenspiel mit Tower-Defense-Mechaniken in Java und libGDX.
Spieleentwicklung ist ein besonders geeignetes Anwendungsgebiet: Spiele verwalten
komplexe, sich ständig ändernde Zustände, viele Objekte interagieren eng
miteinander, und Erweiterbarkeit ist von Anfang an gefordert. Jedes eingesetzte
Pattern löst dabei ein konkretes Architekturproblem, das sich ohne das Muster in
Kopplung, Codeduplikation oder mangelnder Wartbarkeit niedergeschlagen hätte.

\subsection{Spielkonzept}

Der Spieler verteidigt eine Basis gegen zwölf Gegnerwellen, indem er auf einem
20$\times$20-Rasterfeld strategisch Türme platziert. Das Besondere liegt im
kartenbasierten Ressourcensystem: Statt Türme direkt zu kaufen, zieht der Spieler
zu Beginn jeder Runde drei Karten aus einem gemischten Deck. Diese ermöglichen
unterschiedliche Aktionen, Turmbau, Attribut-Buffs, Ressourcengewinn oder
Sondereffekte wie das Verlangsamen von Gegnern.

Jeder Spielzug folgt einem festen Viererzyklus: \textit{Draw} (Karten ziehen,
Energie auffüllen) $\rightarrow$ \textit{Play} (Karten ausspielen, Türme
positionieren) $\rightarrow$ \textit{Wave} (Gegnerwelle läuft, Türme greifen
automatisch an) $\rightarrow$ \textit{Resolve} (kurze Pause, nächster Zug
beginnt). Das Spiel endet mit Sieg nach zwölf abgewehrten Wellen oder mit
Niederlage, sobald die Basis auf null Lebenspunkte fällt. 